\section{Discussion}
\label{section:35-GAI-Discussion}
The proposed \gai benchmarking framework is broadly applicable to \gai problems ranging from simple Graph Manipulations through Searching, Planning and Scheduling up to predictive tasks on graphs. 
The defining features of the benchmark are the focus on problems rather than tasks, and the experimental approach to trade-off efficiency, execution and cost in evaluation. 

Following the \gai benchmark methodology ensures that a collection of problems is built up over time, and can be extended creating a `library' of problems at varying levels of difficulty that can evaluate systems. 

Next, Chapter~\ref{chapter:50-Community-Detection-Case-Study} will step through an example of benchmark development using the domain problem of \textit{Population Segmentation}, which aligns with the \gai problem(s) of \textit{Community Detection, Graph Clustering} and \textit{Graph Partitioning}.